% the page size

\documentclass{article}%
\usepackage{hyperref}
\usepackage{color}
\usepackage{amsmath}
\usepackage{amsfonts}
\usepackage{amssymb}
\usepackage{graphicx}
\usepackage{geometry}%
\setcounter{MaxMatrixCols}{30}
%TCIDATA{OutputFilter=latex2.dll}
%TCIDATA{Version=5.50.0.2953}
%TCIDATA{CSTFile=40 LaTeX article.cst}
%TCIDATA{Created=Thursday, October 06, 2016 16:43:40}
%TCIDATA{LastRevised=Monday, September 25, 2017 17:45:39}
%TCIDATA{<META NAME="GraphicsSave" CONTENT="32">}
%TCIDATA{<META NAME="SaveForMode" CONTENT="1">}
%TCIDATA{BibliographyScheme=Manual}
%TCIDATA{<META NAME="DocumentShell" CONTENT="Standard LaTeX\Blank - Standard LaTeX Article">}
%TCIDATA{Language=American English}
%BeginMSIPreambleData
\providecommand{\U}[1]{\protect\rule{.1in}{.1in}}
%EndMSIPreambleData
\geometry{a4paper, top=1cm, left=1cm, right=1cm, bottom=1.0cm, includehead, includefoot}
\graphicspath{{C:/Users/Gauthier/Dropbox/TD/2015_DSGE/chap3/codes/}}
\newtheorem{theorem}{Theorem}
\newtheorem{acknowledgement}[theorem]{Acknowledgement}
\newtheorem{algorithm}[theorem]{Algorithm}
\newtheorem{axiom}[theorem]{Axiom}
\newtheorem{case}[theorem]{Case}
\newtheorem{claim}[theorem]{Claim}
\newtheorem{conclusion}[theorem]{Conclusion}
\newtheorem{condition}[theorem]{Condition}
\newtheorem{conjecture}[theorem]{Conjecture}
\newtheorem{corollary}[theorem]{Corollary}
\newtheorem{criterion}[theorem]{Criterion}
\newtheorem{definition}[theorem]{Definition}
\newtheorem{example}[theorem]{Example}
\newtheorem{exercise}[theorem]{Exercise}
\newtheorem{lemma}[theorem]{Lemma}
\newtheorem{notation}[theorem]{Notation}
\newtheorem{problem}[theorem]{Problem}
\newtheorem{proposition}[theorem]{Proposition}
\newtheorem{remark}[theorem]{Remark}
\newtheorem{solution}[theorem]{Solution}
\newtheorem{summary}[theorem]{Summary}
\newenvironment{proof}[1][Proof]{\noindent\textbf{#1.} }{\ \rule{0.5em}{0.5em}}
\begin{document}

\title{A note on investment costs and the value of capital goods}
\author{G.\ Vermandel}
\maketitle

\section{Model}

It is possible to derive from the standard RBC\ model the endogenous asset
price denoted $Q_{t}=\lambda_{t}^{k}/\lambda_{t}^{c}$, defined as the ratio
between the marginal cost of raising new capitals (Lagrangian multiplier on
capital accumulation constraint $\lambda_{t}^{k}$)\ and the marginal utility
of consumption (Lagrangian multiplier on budget constraint $\lambda_{t}^{c}$).
Variable $Q_{t}$\ also refers to the well-known Tobin's Q published in Tobin
(1969).\ Tobin's Q is the ratio between a physical asset's market value and
its replacement value, this ratio has considerable macroeconomic implications
as it constitutes the standard bridge between financial markets and markets
for goods and services. In steady state, the ratio must be normalised to one
to be consistent with asset pricing theory, such that $\bar{Q}=1$.

To create a wedge between the price of goods and the price of capital goods,
we introduce a congestion cost on the accumulation process of investment goods
as in Christiano and al.\ (2005):%
\begin{equation}
K_{jt}=\varepsilon_{t}^{I}I_{jt}\left(  1-AC_{jt}^{I}\right)  +\left(
1-\delta\right)  K_{jt-1}%
\end{equation}
where $AC_{jt}^{I}$ is an investment cost function such that:%
\begin{align}
AC_{jt}^{I}  &  =\frac{\chi^{I}}{2}\left(  \frac{I_{jt}}{I_{jt-1}}-1\right)
^{2}\\
&  =f\left(  \frac{I_{jt}}{I_{jt-1}}\right)
\end{align}
This cost function has no implications on the steady state of the model, with
$\overline{AC}^{I}=0$.

\subsection{Households}

The Lagrangian problem is given by:%

\begin{align*}
\mathcal{L}_{t}  &  =\mathbb{E}_{t}%
%TCIMACRO{\dsum \limits_{\tau=0}^{\infty}}%
%BeginExpansion
{\displaystyle\sum\limits_{\tau=0}^{\infty}}
%EndExpansion
\varepsilon_{t+\tau}^{C}\beta^{\tau}\{\frac{C_{jt+\tau}^{1-\sigma_{C}}%
}{1-\sigma_{C}}-\chi\frac{H_{jt+\tau}{}^{1+\sigma_{L}}}{1+\sigma_{L}}\\
&  +\lambda_{t+\tau}^{c}\left[  W_{t+\tau}H_{jt+\tau}+Z_{t+\tau}K_{jt-1+\tau
}+R_{t-1+\tau}B_{jt-1+\tau}-C_{jt+\tau}-I_{jt+\tau}-B_{jt+\tau}-T_{jt+\tau
}\right] \\
&  +\lambda_{t+\tau}^{k}\left[  \varepsilon_{t+\tau}^{I}I_{jt+\tau}\left(
1-\textcolor{red}{f\left(\frac{I_{jt+\tau}}{I_{jt-1+\tau}}\right)}\right)
+\left(  1-\delta\right)  K_{jt-1+\tau}-K_{jt+\tau}\right]  \}\\
&  -
\end{align*}


If you're not comfortable with infinite sums, you can develop the sum for
$\tau=0$ and $\tau=1$ (we don't mind for value of tau above 1):%

\begin{align*}
\mathcal{L}_{t}  &  =\varepsilon_{t}^{C}\left(  \frac{C_{jt}^{1-\sigma_{C}}%
}{1-\sigma_{C}}-\chi\frac{H_{jt}{}^{1+\sigma_{L}}}{1+\sigma_{L}}\right) \\
&  +\varepsilon_{t}^{C}\lambda_{t}^{c}\left[  W_{t}H_{jt}+Z_{t}K_{jt-1}%
+R_{t-1}B_{jt-1}-C_{jt}-I_{jt}-B_{jt}-T_{jt}\right] \\
&  +\varepsilon_{t}^{C}\lambda_{t}^{k}\left[  \varepsilon_{t}^{I}I_{jt}\left(
1-\textcolor{red}{f\left(\frac{I_{jt}}{I_{jt-1}}\right)}\right)  +\left(
1-\delta\right)  K_{jt-1}-K_{jt}\right] \\
&  +\varepsilon_{t+1}^{C}\beta\left(  \frac{C_{jt+1}^{1-\sigma_{C}}}%
{1-\sigma_{C}}-\chi\frac{H_{jt+1}{}^{1+\sigma_{L}}}{1+\sigma_{L}}\right) \\
&  +\varepsilon_{t+1}^{C}\lambda_{t+1}^{c}\beta\left[  W_{t+1}H_{jt+1}%
+Z_{t+1}K_{jt}+R_{t}B_{jt}-C_{jt+1}-I_{jt+1}-B_{jt+1}-T_{jt+1}\right] \\
&  +\varepsilon_{t+1}^{C}\lambda_{t+1}^{k}\beta\left[  \varepsilon_{t+1}%
^{I}I_{jt+1}\left(
1-\textcolor{red}{f\left(\frac{I_{jt+1}}{I_{jt}}\right)}\right)  +\left(
1-\delta\right)  K_{jt}-K_{jt+1}\right] \\
&  +\varepsilon_{t+2}^{C}\beta^{2}\left(  \frac{C_{jt+2}^{1-\sigma_{C}}%
}{1-\sigma_{C}}-\chi\frac{H_{jt+2}{}^{1+\sigma_{L}}}{1+\sigma_{L}}\right) \\
&  +....
\end{align*}


Optimal condition on investment:
\[
-\varepsilon_{t}^{C}\lambda_{t}^{c}+\varepsilon_{t}^{C}\lambda_{t}%
^{k}\varepsilon_{t}^{I}\frac{\partial\left[  I_{jt}\left(  1-AC_{jt}%
^{I}\right)  \right]  }{\partial I_{jt}}-\varepsilon_{t+1}^{C}\lambda
_{t+1}^{k}\beta\varepsilon_{t+1}^{I}\frac{\partial\left[  I_{jt+1}%
AC_{jt+1}^{I}\right]  }{\partial I_{jt}}=0
\]


Only one first order conditions is different from what we had in the standard
RBC model:%
\begin{align*}
\left(  \partial C_{jt}\right)   &  :\lambda_{t}^{c}=C_{jt}{}^{-\sigma_{C}}\\
\left(  \partial H_{jt}\right)   &  :\chi H_{jt}^{\sigma^{L}}=\lambda_{t}%
^{c}W_{t}\\
\left(  \partial B_{jt}\right)   &  :\mathbb{E}_{t}\left\{  \frac
{\varepsilon_{t}^{C}\lambda_{t}^{c}}{\varepsilon_{t+1}^{C}\lambda_{t+1}^{c}%
}\right\}  =\beta R_{t}\\
\left(  \partial I_{jt}\right)   &  :\frac{\lambda_{t}^{k}}{\lambda_{t}^{c}%
}\varepsilon_{t}^{I}\frac{\partial\left[  I_{jt}\left(  1-AC_{jt}^{I}\right)
\right]  }{\partial I_{jt}}-\beta\frac{\varepsilon_{t+1}^{C}\lambda_{t+1}^{k}%
}{\varepsilon_{t}^{C}\lambda_{t}^{c}}\varepsilon_{t+1}^{I}\frac{\partial
\left[  I_{jt+1}AC_{jt+1}^{I}\right]  }{\partial I_{jt}}=1\\
\left(  \partial K_{jt}\right)   &  :\varepsilon_{t}^{C}\lambda_{t}^{k}%
=\beta\mathbb{E}_{t}\left\{  \varepsilon_{t+1}^{C}\left[  \overset{}{\left(
1-\delta\right)  }\lambda_{t+1}^{k}+\lambda_{t+1}^{c}Z_{t+1}\right]  \right\}
\end{align*}
We define the asset price $q_{t}$ by the relative ratio of the marginal
utility of buying capital goods $\lambda_{t}^{k}$ over the marginal utility of
consumption $\lambda_{t}^{c}$ with $q_{t}=\lambda_{t}^{k}/\lambda_{t}^{c}$.
When $q_{t}>1$, households are willing to invest in physical capital rather
than consuming, this can be interpreted as an increase in asset prices with
respect to its steady state. We can thus replace $\lambda_{t}^{k}$\ by the
asset price equation $\lambda_{t}^{k}=q_{t}\lambda_{t}^{c}$. The substitution
affects first order conditions on $\left(  \partial I_{jt}\right)  $ and
$\left(  \partial K_{jt}\right)  $:%
\begin{align*}
\left(  \partial I_{jt}\right)   &  :\frac{\lambda_{t}^{k}}{\lambda_{t}^{c}%
}\varepsilon_{t}^{I}\frac{\partial\left[  I_{jt}\left(  1-AC_{jt}^{I}\right)
\right]  }{\partial I_{jt}}-\beta\frac{\varepsilon_{t+1}^{C}\lambda_{t+1}^{k}%
}{\varepsilon_{t}^{C}\lambda_{t}^{c}}\varepsilon_{t+1}^{I}\frac{\partial
\left[  I_{jt+1}AC_{jt+1}^{I}\right]  }{\partial I_{jt}}=1\\
&  :\frac{q_{t}\lambda_{t}^{c}}{\lambda_{t}^{c}}\varepsilon_{t}^{I}%
\frac{\partial\left[  I_{jt}\left(  1-AC_{jt}^{I}\right)  \right]  }{\partial
I_{jt}}-\beta\frac{\varepsilon_{t+1}^{C}q_{t+1}\lambda_{t+1}^{c}}%
{\varepsilon_{t}^{C}\lambda_{t}^{c}}\varepsilon_{t+1}^{I}\frac{\partial\left[
I_{jt+1}AC_{jt+1}^{I}\right]  }{\partial I_{jt}}=1\\
&  :q_{t}\varepsilon_{t}^{I}\frac{\partial\left[  I_{jt}\left(  1-AC_{jt}%
^{I}\right)  \right]  }{\partial I_{jt}}-\frac{1}{R_{t}}\mathbb{E}_{t}%
q_{t+1}\varepsilon_{t+1}^{I}\frac{\partial\left[  I_{jt+1}AC_{jt+1}%
^{I}\right]  }{\partial I_{jt}}=1
\end{align*}
and for $\left(  \partial K_{jt}\right)  $:%
\begin{align*}
\left(  \partial K_{jt}\right)   &  :\varepsilon_{t}^{C}q_{t}\lambda_{t}%
^{c}=\beta\mathbb{E}_{t}\left\{  \varepsilon_{t+1}^{C}\left[  \overset
{}{\left(  1-\delta\right)  }q_{t+1}\lambda_{t+1}^{c}+\lambda_{t+1}^{c}%
Z_{t+1}\right]  \right\} \\
&  :q_{t}=\beta\mathbb{E}_{t}\left\{  \frac{\lambda_{t+1}^{c}\varepsilon
_{t+1}^{C}}{\varepsilon_{t}^{C}\lambda_{t}^{c}}\left[  \overset{}{\left(
1-\delta\right)  }q_{t+1}+Z_{t+1}\right]  \right\} \\
&  :q_{t}=\mathbb{E}_{t}\left\{  \frac{1}{R_{t}}\left[  \overset{}{\left(
1-\delta\right)  }q_{t+1}+Z_{t+1}\right]  \right\} \\
&  :R_{t}=\frac{\overset{}{\left(  1-\delta\right)  \mathbb{E}_{t}}%
q_{t+1}+\mathbb{E}_{t}Z_{t+1}}{q_{t}}%
\end{align*}
these two equations are really popular in DSGE\ modelling.\ Many papers
re-employ this investment cost function as it generate a dynamic of investment
consistent with empirical evidence and data friendly (response of invesment is
hump-shaped as in VARs models).

\section{Model summary}

We have to include one new variables $\{\textcolor{red}{q_{t}}\}$, and one
corresponding equation to obtain a determinate system of equations:%

\begin{align*}
&  :\beta R_{t}=\mathbb{E}_{t}\left\{  \frac{\varepsilon_{t}^{C}}%
{\varepsilon_{t+1}^{C}}\left(  \frac{C_{jt+1}}{C_{jt}}\right)  ^{\sigma_{C}%
}\right\} \\
&  :W_{t}=\chi H_{t}^{\sigma_{L}}C_{t}^{\sigma_{C}}\\
&  :R_{t}=\frac{\mathbb{E}_{t}\left\{  \left(  1-\delta\right)
\textcolor{red}{q_{t+1}}+Z_{t+1}\right\}  }{\textcolor{red}{q_{t}}}\\
&  :K_{t}=\varepsilon_{t}^{I}\left(
1-\textcolor{red}{f\left(\frac{I_{jt}}{I_{jt-1}}\right)}\right)  I_{t}+\left(
1-\delta\right)  K_{t-1}\\
&
:\textcolor{red}{ q_{t}\varepsilon_{t}^{I}\frac{\partial\left[ I_{jt}\left( 1-AC_{jt}^{I}\right) \right] }{\partial I_{jt}}-\frac{1}{R_{t}}\mathbb{E}_{t}q_{t+1}\varepsilon_{t+1}^{I}\frac{\partial\left[ I_{jt+1}AC_{jt+1}^{I}\right] }{\partial I_{jt}}=1 }\\
&  :Y_{t}=\varepsilon_{t}^{A}K_{t-1}^{\alpha}H_{t}{}^{1-\alpha},\\
&  :Z_{t}=\alpha\frac{Y_{t}}{K_{t-1}}\\
&  :W_{t}=\left(  1-\alpha\right)  \frac{Y_{t}}{H_{t}}\\
&  :Y_{t}=C_{t}+I_{t}+g^{y}\bar{Y}\varepsilon_{t}^{G}\\
&  :\log\left(  \varepsilon_{t}^{A}\right)  =\rho_{A}\log\left(
\varepsilon_{t-1}^{A}\right)  +\eta_{t}^{A}\\
&  :\log\left(  \varepsilon_{t}^{G}\right)  =\rho_{G}\log\left(
\varepsilon_{t-1}^{G}\right)  +\eta_{t}^{G}\\
&  :\log\left(  \varepsilon_{t}^{C}\right)  =\rho_{C}\log\left(
\varepsilon_{t-1}^{C}\right)  +\eta_{t}^{C}\\
&  :\log\left(  \varepsilon_{t}^{I}\right)  =\rho_{I}\log\left(
\varepsilon_{t-1}^{I}\right)  +\eta_{t}^{I}%
\end{align*}


\section{Steady state}

The asset price equation is normalized to unity in steady state:%
\[
\bar{q}=1
\]


\section{Dynare writing}

\begin{enumerate}
\item Declare one new endogenous variable $\{\textcolor{red}{q_{t}}\}$

\item Declare the investment cost $\chi$\ as a parameter and calibrate it to
$\chi$\ =4

\item Include the one dynamic equation defining $\{\textcolor{red}{q_{t}}\}$
and modify equations affected by capital utilization.

\item Define the steady state of $\{\textcolor{red}{q_{t}}\}$\ in the init block.
\end{enumerate}

\section{Tips}

First, regarding investment costs derivation:%
\begin{align*}
\frac{\partial\left[  I_{t}AC_{t}^{I}\right]  }{\partial I_{t}}  &
=\frac{\chi}{2}\left(  3\left(  \frac{I_{t}}{I_{t-1}}\right)  ^{2}%
+1-4\frac{I_{t}}{I_{t-1}}\right)  ,\\
\frac{\partial\left[  I_{t+1}AC_{t+1}^{I}\right]  }{\partial I_{t}}  &
=\chi\left(  \left(  \frac{I_{t+1}}{I_{t}}\right)  ^{2}-\left(  \frac{I_{t+1}%
}{I_{t+1}}\right)  ^{3}\right)  .
\end{align*}


Second you may have an headache reading this equation:
\[
:q_{t}\varepsilon_{t}^{I}\frac{\partial\left[  I_{jt}\left(  1-AC_{jt}%
^{I}\right)  \right]  }{\partial I_{jt}}-\frac{1}{R_{t}}\mathbb{E}_{t}%
q_{t+1}\varepsilon_{t+1}^{I}\frac{\partial\left[  I_{jt+1}AC_{jt+1}%
^{I}\right]  }{\partial I_{jt}}=1
\]
however, up to first order, this equation boils down to:%
\[
\hat{q}_{t}=\chi(\hat{\imath}_{t}-\hat{\imath}_{t-1})-\beta\chi(\hat{\imath
}_{t+1}-\hat{\imath}_{t})
\]

\end{document}